\chapter{Conclusions and Future Work}\label{chapter:chap9}

\section{Summary of Contributions}\label{sect:contributions}

The thesis aims to design, implement and evaluate a security analysis framework for Rust's ecosystem. The results demonstrate that automated multi-dimensional analysis is effective and feasible, with some limitations including false positives and the need for manual verification.

\begin{enumerate}[leftmargin=*,itemsep=2pt]
\item \textbf{A modular and scalable analysis framework.} We designed and implemented a producer consumer architecture with a Redis-based work queue that enforces topological ordering. The system supports worker containers, handles crash recovery and graceful shutdown. A full ecosystem scan of 226,546 crates can be completed with consumer hardware. 

\item \textbf{Five integrated modules.} The framework combines:
\begin{itemize}[leftmargin=*,itemsep=1pt]
\item Vulnerability scanning using \texttt{cargo-audit}, with a CVSS scoring implementation
\item Secret detection from \texttt{gitleaks} with rule pattern matching and entropy analysis
\item Heuristic \texttt{build.rs} analysis covering multiple patterns (network calls, process spawning, raw IPs, link directives, free TLDs, entropy)
\item Multi-algorithm typosquatting detection with weighted combined scoring
\item LLM-based malicious code scoring via LM Studio's API
\end{itemize}


\item \textbf{The largest empirical security assessment of crates.io to date.} Scanning 226,546 crates (85.6\% of the registry) resulted in the following key findings:
\begin{itemize}[leftmargin=*,itemsep=1pt]
\item 32.1\% of scanned crates (72,773) have known vulnerabilities; 82.6\% of findings are informational severity, while $\sim$11\% are High/Critical.
\item 882,955 secrets were detected across 18,571 crates, but 95.8\% are \texttt{generic-api-key} false positives.
\item 1.41\% of scanned crates (3,194) spawn processes in \texttt{build.rs}, representing a compile-time attack surface present across thousands of crates.
\item 38,495 name-similarity pairs were detected; 891 with very high similarity ($\ge$95). Some have high download asymmetry suggesting possible impersonation.
\item 13\% of crates (34,366) could not be analyzed, a significant blind spot reflecting private repositories and limitations of the current implementation, missing \texttt{Cargo.lock} files, and clone failures.
\end{itemize}

\item \textbf{Evidence-based risk characterization.} Beyond raw counts, the analysis revealed structural patterns with practical implications:
\begin{itemize}[leftmargin=*,itemsep=1pt]
\item Popular crates are safer high download crates have a 25.2\% vulnerability rate versus 32.1\% for medium download crates.
\item Newest crates are most vulnerable (30.9\% rate for crates under 6 months), likely due to unpatched dependencies and less review.
\item Crates shipping executables have a 67.7\% vulnerability rate; more than double the average.
\item The Substrate/Polkadot ecosystem is the largest source of transitive vulnerability risk due to a single advisory from the dependency \texttt{sp-core}.
\end{itemize}

\item \textbf{A negative result for LLM-based module.} Deploying a 14B-parameter LLM at ecosystem scale was not possible due to hardware limitations, scanning the most vulnerable crates didn't produce meaningful findings.

\end{enumerate}

\section{Future Directions}\label{sect:future-directions}

The framework and findings show several opportunities for future work:

\begin{enumerate}[leftmargin=*,itemsep=2pt]
\item \textbf{Version-level dependency analysis.} The current implementation tracks vulnerabilities at the crate level, without tracking specific versions. The \texttt{versions} and \texttt{dependencies} tables already contain version-level data, with the pipeline logic being the missing component.

\item \textbf{Integration with crates.io infrastructure.} The \texttt{build.rs} heuristic flags could trigger automated warnings on the crate's crates.io page, alerting users to compile-time execution before they add the dependency.

\item \textbf{Gitleaks rule set specialization.} 
Gitleaks has a very high false positive rate in this context. Even after removing test and example paths the remaining are too many for manual review.

\item \textbf{Caching and incremental scanning.} Caching cargo audit's artifacts such as different compiler versions and implementing incremental analysis that only re-analyzes crates whose source code or dependency graph has changed would reduce the scan time.

\item \textbf{Sandboxed build script execution.} A very important gap in the current implementation is the lack of understanding what the build.rs script executes at runtime. Implementing a sandbox and monitoring behavior would be a great addition, especially with a classifier to determine possible maliciousness.

\item \textbf{Continuous monitoring.} Integrating the framework with crates.io's real-time feed (instead of the current "periodic" database dump) would enable analysis of newly published crates. The redis queue already is capable of handling a continuous addition of crates.

\item \textbf{Improved LLM integration.} Deploying a larger model with RAG for cross-file context, would be a great addition, especially deployed on a cluster for fast inference, with the current implementation was not properly tested due to hardware limitations.

\item \textbf{Metadata anomaly detection.} Analyzing patterns in maintainer and ownership changes, release frequency, and version-tag rewriting would complement the current analysis methods.

\end{enumerate}

